\documentclass[11pt]{article}
\usepackage[margin=1in]{geometry}
\usepackage{amsmath,amssymb,amsthm,mathtools}
\usepackage[hidelinks]{hyperref}
\newtheorem{theorem}{Theorem}
\newtheorem{lemma}[theorem]{Lemma}
\newtheorem{corollary}[theorem]{Corollary}
\newcommand{\C}{\mathbb C}
\newcommand{\Q}{\mathbb Q}
\newcommand{\E}{\mathbb E}
\newcommand{\fall}[2]{#1^{\underline{#2}}}
\title{Uniform Gamma and simplex monomial windows, with every polynomial marker}
\author{Internal proof and exact controls; no priority claim}
\date{16 November 2026}
\begin{document}
\maketitle
\noindent\textbf{Progress and boundary.}
The preceding \href{MIXED_PRIME_RESIDUAL_WINDOWS.tex}{mixed-prime packet}
is revalidated. Its unresolved first residual branches are retained.
Here an elementary differential equation supplies a different, uniform
conclusion: for every exponent, every positive real Gamma/Dirichlet shape,
and every starting index, a degree-dependent moment window is nonzero.
Every polynomial marker is either identically invisible or has an explicit
window. This proves SFC for the expanded block-sum classes below, including
mixed exponents. For the two-term polynomial $a+bY^6$ the unmarked bound
is seven, not the conjectured two. Full exponent-six SFC, unrestricted SFC
and general FC in dimension at least three remain open.
The argument is derived here from literal moment coefficients; it needs no
new literature input and makes no novelty, external-review or formal-proof claim.

\section{Exact marker polynomials}
Let $d_i\geq0$ be integers, $D=\sum_i d_i\geq1$, and $\alpha_i>0$ real.
Put $Z=\prod_i Y_i^{d_i}$ for independent $Y_i\sim\operatorname{Gamma}(\alpha_i,1)$,
or $Z=\prod_i T_i^{d_i}$ for $T\sim\operatorname{Dir}(\alpha)$.
The Dirichlet law has at least two coordinates. In either case $Z$ is
nonconstant. Write $A=\sum_i\alpha_i$, let $s$ be the number of positive
$d_i$, and put $A_{\rm inact}=\sum_{d_i=0}\alpha_i$.
For $F=a+bZ$, $a,b\in\C$, assume $b\ne0$. Denote its marked moments by
$\mu_n=\E[M F^n]$, including $\mu_0=\E[M]$.

For Gamma variables, write $M=\sum_e m_eY^e$ and define
\[
 R(t)=\sum_e m_e\prod_i(\alpha_i+d_it)_{e_i},\qquad
 c_k=\prod_i(\alpha_i)_{d_i k}.                         \tag{1}
\]
For Dirichlet variables, homogenize the marker on $\sum T_i=1$ to any
degree $E\geq\deg M$:
\[
 \widetilde M(T)=\sum_e m_eT^e(\sum T_i)^{E-|e|}
                 =\sum_{|e|=E}\widetilde m_e T^e.
\]
In this case set
\[
 R(t)=\sum_{|e|=E}\widetilde m_e\prod_i(\alpha_i+d_it)_{e_i},\qquad
 c_k=\frac{\prod_i(\alpha_i)_{d_i k}}{(A)_{Dk+E}}.       \tag{2}
\]
Both formulas give $\E[M Z^k]=c_kR(k)$ exactly. All $c_k$ are strictly
positive. If $R\ne0$, put $e=\deg R$; in (2), $e\leq E$.
Inactive coordinates and cancellations in $R$ are retained.

\begin{theorem}[All markers, all starts]\label{thm:window}
For either law and every complex $a,b$ with $b\ne0$:
\begin{enumerate}
\item $R=0$ if and only if every marked moment vanishes, including $n=0$.
\item If $R\ne0$, every $D+e+1$ consecutive positive marked moments
      contain a nonzero value.
\end{enumerate}
In particular every nonzero $a+bZ$ has a nonzero unmarked moment in every
$D+1$ consecutive positive indices, for all positive real shapes.
For a monomial marker, $e$ is the total marker exponent on active coordinates.
\end{theorem}
For $b=0$, the separate answer is $\mu_n=a^n\E[M]$; a zero-mass marker
can vanish identically for a nonzero constant phase even when $R\ne0$.
The zero phase and the one-coordinate Dirichlet point law are not covered
by the nonconstant assertion.

\section{The differential equation and the simplex singularity}
Put
\[
 Q(t)=\prod_{i:d_i>0}\prod_{j=0}^{d_i-1}(\alpha_i+d_it+j),\qquad
 q_D=\prod_{i:d_i>0}d_i^{d_i}>0.
\]
For Gamma let $S(t)=1$; for Dirichlet let
$S(t)=(A+E+Dt)_D$ and $s_D=D^D$.
Then $c_{k+1}/c_k=Q(k)/S(k)$ in the respective formulas above.
For an integer $N\geq0$ write
\[
 g_N(x)=\sum_{k=0}^N\binom Nk c_k R(k)x^k,
 \qquad \theta=x\frac{d}{dx}.
\]
Coefficient comparison, including both endpoints, gives
\[
 \big[\theta S(\theta-1)R(\theta-1)
       -x(N-\theta)Q(\theta)R(\theta+1)\big]g_N=0.    \tag{3}
\]
No commutation of $x$ past $\theta$ is made: the displayed $x$ acts last.
Indeed the coefficient of $x^k$ on both sides before subtraction is
$kR(k-1)S(k-1)\binom Nk c_kR(k)$ for $k\geq1$.
If $N\geq e$ then $g_N\ne0$, since a nonzero degree-$e$ polynomial
cannot vanish at all integers $0,\ldots,N$.

\begin{lemma}\label{lem:mult}
For $R\ne0$ and $N\geq e$, every nonzero root of $g_N$ has multiplicity
at most $D+e$.
\end{lemma}
\begin{proof}
Put $L=D+e+1$ and let $r_e$ be the leading coefficient of $R$.
For Gamma the leading ordinary derivative coefficient of (3) is
$q_Dr_e x^{L+1}$, nonzero for $x\ne0$. A root of order $t\geq L$
would give a unique nonzero term of order $t-L$ in (3), a contradiction.

For Dirichlet the coefficient is $r_e x^L(s_D+q_Dx)$.
There is one extra point $x_*=-s_D/q_D$. We check it rather than treating
it as ordinary. If $P(\theta)-xB(\theta)$ denotes (3), then
\begin{align*}
 P(t)&=tS(t-1)R(t-1),& B(t)&=(N-t)Q(t)R(t+1),\\
 \frac{P_{L-1}}{P_L}
   &=A+E-\frac{D+1}{2}+\rho-e,&
 \frac{B_{L-1}}{B_L}
   &=-N+\sum_{d_i>0}\alpha_i+\frac{D-s}{2}+\rho+e,
 \end{align*}
where $\rho=r_{e-1}/r_e$, interpreted as zero if $e=0$.
The identity
$\theta^L=x^L\partial^L+\binom L2 x^{L-1}\partial^{L-1}+\cdots$
shows that a root of order $t\geq L-1$ at $x_*$ must satisfy
\[
 t=L-1+\frac{P_{L-1}}{P_L}-\frac{B_{L-1}}{B_L}
   =N+\beta,\qquad
 \beta=A_{\rm inact}+E-e+\frac{s-1}{2}>0.             \tag{4}
\]
Explicitly, compare the term of order $t-L+1$ from the leading two
ordinary derivatives, using the simple zero of their leading coefficient.
All lower derivatives have higher order. Strict positivity follows from
$E\geq e$: if $s\geq2$ its last term is positive; if $s=1$, the
at-least-two-coordinate hypothesis gives $A_{\rm inact}>0$.
But $\deg g_N\leq N$, so (4) is impossible. Thus the extra singularity
does not invalidate the bound (and actually permits order at most $L-2$).
\end{proof}

\begin{proof}[Proof of Theorem~\ref{thm:window}]
If $a\ne0$, then $\mu_j=a^j g_j(b/a)$. For every coefficient sequence,
\[
 \fall Nj\,g_{N-j}
    =\prod_{i=0}^{j-1}(N-i-\theta)g_N.               \tag{5}
\]
Consequently $L$ consecutive zero moments ending at $N\geq L$ imply
$g_N$ and its first $L-1$ derivatives vanish at $b/a\ne0$.
This contradicts Lemma~\ref{lem:mult}; $N\geq L>e$ ensures $g_N\ne0$.
If $a=0$, then $\mu_j=b^j c_jR(j)$; $L>e$ distinct positive integer
zeros would force $R=0$. Conversely $R=0$ makes every moment zero by
binomial expansion. The nonzero-window assertion gives the other
direction, as well as exclusion of every eventual zero tail.
\end{proof}

\section{Ordinary coefficient ideals over every rational algebra}
Here use rational positive shapes and the unmarked law. Define the
homogeneous polynomials
$\mu_j(U,V)=\E[(U+VZ)^j]\in\Q[U,V]$.

\begin{corollary}\label{cor:ideal}
For any start $n\geq1$, put $N=n+D$ and
$I=(\mu_n,\ldots,\mu_N)$. This ordinary ideal contains a homogeneous
regular sequence of degrees $n,N$. In particular
\[
 (U,V)^{2n+D-1}\subseteq I,\qquad
 \dim_\Q\Q[U,V]/I\leq n(n+D).                       \tag{6}
\]
Over every $\Q$-algebra, zero moments in that window force the actual
phase coefficient ideal $(a,b)$ to be nilpotent at exponent $2n+D-1$.
Thus the phase itself and every polynomially marked power vanish at all
indices at least that exponent. No sharpness of these bounds for $I$
is asserted.
\end{corollary}
\begin{proof}
Theorem~\ref{thm:window} rules out projective common zeros; at $U=0$
or $V=0$ each $\mu_j$ is also nonzero. For a zero of $\mu_n$ some
$\mu_{n+j}$, $1\leq j\leq D$, is nonzero. Thus
\[
 G_t=\sum_{j=1}^D t^{j-1}U^{D-j}\mu_{n+j}
\]
has at most $D-1$ bad values of $t$ at that point. There are at most $n$
distinct projective roots of $\mu_n$. Some integer
$0\leq t\leq n(D-1)$ therefore makes $\mu_n,G_t$ coprime.
They form a binary complete intersection of degrees $n,N$, with Hilbert
series $(1+\cdots+z^{n-1})(1+\cdots+z^{N-1})$, length $nN$ and top
degree $n+N-2$. Its containment in $I$ proves (6). These are statements
about ordinary ideals, so base change to arbitrary coefficient rings is
valid. For Dirichlet, eliminate one simplex coordinate: the nonconstant
rational polynomial $Z$ has a nonzero nonconstant coefficient, a rational
unit, so the actual phase coefficient ideal is still $(a,b)$.
\end{proof}

\section{Strong Factorial classes from actual expanded support}
Partition factorial variables into disjoint nonempty blocks of sizes $r_i$,
put $L_i=\sum_{j\in B_i}Y_j$, and let $d_i\geq1$.
Then $L_i$ are independent Gamma variables with shape $r_i$.
For $a,b\ne0$ the actual expanded polynomial
\[
 F=a+b\prod_i L_i^{d_i}
\]
has exactly $C=1+\prod_i\binom{r_i+d_i-1}{d_i}$ monomials.
\begin{corollary}\label{cor:sfc}
If $C\geq D+1$, this class satisfies SFC at every starting index.
In particular it does so for arbitrary exponents when every $r_i\geq2$.
Disjoint independent linear adjunctions preserve the assertion as in the
\href{SFC_WINDOWS_AND_ROCKS_AUDIT.tex}{checked Rocks construction}.
The cases $a=0$ or $b=0$ have nonzero pure moments directly.
\end{corollary}
\begin{proof}
Theorem~\ref{thm:window} gives a $D+1$ window. If each $r_i\geq2$, then
$\prod_i\binom{r_i+d_i-1}{d_i}\geq\prod_i(d_i+1)\geq D+1$.
For adjunctions, each nonzero new independent linear term adds one to
the window bound and to actual support. Degenerate base phases have
window one, so they obey the same assertion after adjunction.
\end{proof}
For example $a+b(Y_1+Y_2)^6$ has eight actual monomials and a seven-moment
window, for every start. This is a full SFC class. The two-monomial
$a+bY^6$ receives the same seven-moment bound, which is not full SFC.

\section{Original Gaussian markers and executable scope}
Put $Y_i=\frac12\sum_{j\in B_i}X_j^2$ for standard independent real
Gaussians. Then $\alpha_i=|B_i|/2$ is allowed without any integer-shape
restriction. Average any original polynomial marker over the independent
orthogonal block groups using the exact Reynolds formula in
\href{moment_bridges.py}{moment\_bridges.py}; it gives a polynomial
$\overline M(Y)$ with exactly the same moments against every $F^n$.
Apply (1) to this actual averaged marker. Thus $R=0$ decides its whole
zero sequence; otherwise Theorem~\ref{thm:window} gives the stated window.
Odd markers can average to zero and are not misclassified as nonzero.
The actual Gaussian phase coefficient ideal is $(a,b)$: its constant is
$a$ and a pure coordinate monomial has coefficient $b/2^D$. Thus the
ordinary nilpotence bound also transfers without taking radicals.
The unmarked Gaussian polynomial $a+b\prod_i(\frac12\sum X_j^2)^{d_i}$
also satisfies the support-count version of Corollary~\ref{cor:sfc}
(a Gaussian moment analogue of SFC). No independent Gaussian-linear
adjunction is inferred from the exponential-variable argument.

\href{hypergeometric_windows.py}{hypergeometric\_windows.py} accepts
positive rational shapes and exact algebraic marker coefficients, returns
the literal $Q,S,R$ and both parameter singularities, and exports either
an exact zero-sequence identity or the universal window. The proofs above
allow all positive real shapes; finite executable inputs are narrower.
The ordinary-ideal exporter supplies an actual coprime $G_t$ and retains
the distinction between its regular subideal and the entire window ideal.
\href{verify_hypergeometric_windows.py}{Exact controls} compare the
coefficient formulas to independent direct Gamma, Dirichlet and original
Gaussian integration, expand the ordinary differential operator, and
check its singular indicial coefficient, ideal inclusions and boundary cases.
They support the written all-index proofs rather than replace them.
\end{document}
